\subsection{Every higher central lifting fiber and the complete class-three fiber group}

\begin{statement}{E.40}{Proved}{Higher intrinsic fibers have finite matrix presentations; the next one is evaluated}
For the full cubic pro-three group $N$ put
\[
 Q_r=N/\gamma_{r+1}N,\qquad
 L_r=\gamma_rN/\gamma_{r+1}N
       =\operatorname{Lie}_r(V),\qquad r\ge3.
 \tag{E40.1}
\]
The equality is the actual graded commutator isomorphism
over $\mathbb Z_3$. Set
$A_r=\operatorname{Lie}_r(A_\tau)$ and
$C_r=\operatorname{Lie}_r(C_V)$, using the integral
matrices of E.37. Here $c$ denotes complex conjugation;
$L_r^{\langle c\rangle}$ means its fixed sublattice.
Every nonempty geometric-conjugacy fiber of
\[
 \operatorname{Sect}(Q_r\rtimes G_{\mathbb Q})
      \longrightarrow
 \operatorname{Sect}(Q_{r-1}\rtimes G_{\mathbb Q})
 \tag{E40.2}
\]
is a torsor under the explicitly presented compact group
\[
 L_r/(A_r-1)L_r^{\langle c\rangle}
       =\operatorname{coker}\bigl((A_r-1)(1+C_r):L_r\to L_r\bigr).
 \tag{E40.3}
\]
There is no further stabilizer quotient. The central
cocycles acting in these fibers are all unramified away
from three, including when the lower section is ramified.

For the first subsequent step, namely $Q_3\to Q_2$,
every nonempty fiber is exactly a torsor under
\[
 \mathbb Z_3^{166}\ \oplus\
 (\mathbb Z/3)^7\ \oplus\
 (\mathbb Z/9)^{56}\ \oplus\
 (\mathbb Z/27)^{25}\ \oplus\
 (\mathbb Z/81)^3.
 \tag{E40.4}
\]
Thus the nonempty fibers at this stage have a fully
computed integral parameter group, while E.38 supplies
an entire abelian class whose class-two fibers all have
empty preimage at this stage.

The presentation (E40.3) is effectively finite in every
degree. Its free rank is
\[
 \frac{n_r-t_r}{2},\qquad
 n_r=\frac1r\sum_{d\mid r}\mu(d)10^{r/d},\qquad
 t_r=\frac1r\sum_{d\mid r}\mu(d)
       \bigl(\operatorname{tr}(c^d\mid V)\bigr)^{r/d},
 \tag{E40.5}
\]
where the last trace is $10$ for even $d$ and $-2$
for odd $d$. Its torsion is killed by
\[
 3^{r+b_r},\qquad
 b_r=\max_{r\le j\le2r}(1+v_3(j)).
 \tag{E40.6}
\]
Consequently a Hall-basis matrix calculation modulo
$3^{r+b_r+1}$ determines all its elementary divisors.

\emph{Manuscript: \S51.16, folios 733v--736v, scans 172--178,
Part E(c),(d). The comparison is with the full intrinsic
lower-central quotients; the precise further maps from
the metabelian and single-cusp marked quotients are
given below.}
\end{statement}
\begin{proof}
The free pro-three group $N$ has rank ten. The graded
commutator map identifies its lower-central graded
algebra with the free $\mathbb Z_3$-Lie algebra on $V$;
the exact pro-$p$ statement is
\cite[Theorem 2.1, pp.~4--5]{E_KochKukkukLabute}.
In Magnus coordinates this sends the degree-$r$
leading term of a group commutator to the same Lie
bracket in the tensor algebra. It is functorial for
automorphisms, so the arithmetic action on $L_r$ is
precisely $\operatorname{Lie}_r$ of the action on $V$.
This uses the integral graded quotient, including at
$r=3$; it does not use a class-three exponential with
denominator three.

Conjugation by an element of $N$ acts trivially on
$L_r$, since $[N,\gamma_rN]\subset\gamma_{r+1}N$.
Therefore the action on this central kernel after
pulling back along any lower section is the same
displayed arithmetic action. Twisting the lower
section cannot change that module.

At a rational prime $p\ne3$ the original action is
unramified. The eigenvalues on $V\otimes\mathbb Q_3$
have complex absolute values $p^{1/2}$ or $p$:
the actual isogeny $T_D\times E\to J_{E,D}$ in E.24
gives these toric and elliptic factors. The eigenvalues
on $L_r\otimes\mathbb Q_3$, as a subspace of
$V^{\otimes r}$, therefore have absolute values
$p^{j/2}$ with $r\le j\le2r$. For $r\ge3$ none is
$p$ and none is one.

Let $a$ be a continuous $L_r$-valued cocycle. Wild
inertia at $p$ has zero image because it is pro-$p$.
On tame pro-three inertia, the cocycle is a
homomorphism; its value $v$ on a generator satisfies
$Fv=pv$ by the tame relation. The absence of the
eigenvalue $p$ makes $v=0$ in the torsion-free lattice.
Thus $a$ is zero on every inertia group outside
three, and factors through $G_{\mathbb Q,\{3,\infty\}}$.
This is an equality of cocycle groups, so there is no
inverse-limit ramification assumption in this step.

The same weight calculation excludes a geometric
stabilizer of a lower section. If a nonidentity
element of $Q_{r-1}$ were fixed by its section-twisted
Galois action, take its first nonzero lower-central
coordinate. This coordinate lies in some
$L_i$, $1\le i<r$, and is fixed by the original action:
inner twists act trivially on the graded coordinate.
Every Frobenius eigenvalue there has absolute value
strictly greater than one, contradicting the nonzero
fixed coordinate. Hence
\[
 H^0(G_{\mathbb Q},Q_{r-1}^{\mathrm{section}})=1.
 \tag{E40.7}
\]
For two lifts of a fixed lower section, their pointwise
ratio is an $L_r$-valued cocycle. Central conjugation
changes it by a coboundary. Equation (E40.7) means
that no additional conjugation in the lower group
identifies two such cocycle classes. This proves that
the full geometric-conjugacy fiber is a torsor under
$H^1(\mathbb Q,L_r)$.

This cohomology group is computed directly, rather
than left as its name. Over $K$, the unramified group
$P=\langle\tau,\tau'\rangle$ is free and
$c\tau c=\tau'$. As in E.34, its action and every
cocycle factor through $P\rtimes\langle c\rangle$.
Normalize an additive cocycle to be zero at $c$ by
division by two. It is then freely determined by
$x=a(\tau)\in L_r$; its value at $\tau'$ is $C_rx$.
The coboundaries preserving this normalization are
exactly $(A_r-1)y$ with $y\in L_r^{\langle c\rangle}$.
Since two is a unit, the image of $1+C_r$ is exactly
that fixed sublattice. This gives (E40.3) by an actual
evaluation map. Frobenius at seven has no eigenvalue
one, so $A_r-1$ is rationally injective and the free
rank of its cokernel on the fixed sublattice is the
minus-space dimension. Witt's character formula
\cite[Theorem A.1]{E_AkaWitt}, evaluated at the ten
eigenvalues of $c$, gives (E40.5).

We give a uniform integral bound that makes the
finite matrix calculation conclusive. In $V$ consider
the stable sublattice
\[
 V_0=T_3T_D+j(T_3E).
 \tag{E40.8}
\]
Its matrix in the basis of E.37 is
$B=\left(\begin{smallmatrix}3I_2&0\\U&I_8\end{smallmatrix}\right)$,
and
$B\left(\begin{smallmatrix}I_2&0\\-U&3I_8\end{smallmatrix}\right)=3I_{10}$.
Thus $3V\subset V_0\subset V$. By multilinearity
$3^rL_r\subset\operatorname{Lie}_r(V_0)\subset L_r$.
At the unique $K$-prime above two, the elliptic scalar
is $-2$ and the toric scalar is four. A word with
$a$ elliptic letters has scalar
$(-2)^a4^{r-a}=(-2)^{2r-a}$ on the split lattice.
For $j\ge1$, elementary binomial induction gives
$v_3((-2)^j-1)=1+v_3(j)$. Hence $3^{b_r}(F-1)^{-1}$
preserves $\operatorname{Lie}_r(V_0)$, and
\[
 3^{r+b_r}(F-1)^{-1}L_r\subset L_r.
 \tag{E40.9}
\]
The torsion of $H^1(\mathbb Q,L_r)$ is exactly
$H^0(\mathbb Q,L_r\otimes\mathbb Q_3/L_r)$:
this follows from the long exact sequence and
$H^0(\mathbb Q,L_r\otimes\mathbb Q_3)=0$.
An invariant class has $(F-1)v\in L_r$, so (E40.9)
kills it by $3^{r+b_r}$. This proves (E40.6).

Here are full coordinates and the certificate in
degree three. Use the $330$ brackets
\[
 [[x_i,x_j],x_k],\qquad
 0\le j<i\le9,\quad j\le k\le9,
 \tag{E40.10}
\]
ordered lexicographically in $(i,j,k)$, with the $x_i$
in the order of the $V$-basis of E.37. Antisymmetry
and the identity
\[
 [[x_i,x_j],x_k]
   =[[x_i,x_k],x_j]-[[x_j,x_k],x_i]
       \quad(k<j<i)
 \tag{E40.11}
\]
reduce every bracket to these coordinates. They
span; their number is $(10^3-10)/3=330$, the free
rank of $L_3$, so they are a basis. Applying the
two known ten-dimensional matrices to all three
letters and then (E40.11) constructs $A_3,C_3$
by integer operations.

Their cocycle matrix $P_3=(A_3-1)(1+C_3)$ has rank
$164$: the trace of $C_3$ is
$((-2)^3-(-2))/3=-2$, giving plus and minus dimensions
$164$ and $166$. The program
\texttt{e\_lie3\_fibers.py} records invertible modular
row and column transformations $S,T$ and verifies
\[
 SP_3T\equiv
 \operatorname{diag}(1^{\times73},3^{\times7},
 9^{\times56},27^{\times25},81^{\times3},0^{\times166})
       \pmod{729}.
 \tag{E40.12}
\]
Every scaling operation is a unit modulo three, and
the output includes every operation and both complete
transformation matrices. Here $b_3=2$, so (E40.6)
bounds the torsion exponent by $3^5$. Therefore the
modulo-$3^6$ certificate determines the integral
elementary divisors exactly; none can be hidden above
the tested modulus. The independent stacked matrix
$\binom{A_3-1}{C_3-1}$ computes the invariants of
$L_3\otimes\mathbb Q_3/L_3$ and gives the same
torsion factors. The program also verifies the exact
split-lattice identities used in (E40.9), with scalar
values $64,-32,16,-8$. This proves (E40.4).
The exact quotient presentation (E40.3) supplies full
$\mathbb Z_3$ coordinates; (E40.12) gives their Smith
coordinates modulo $729$, and the same elementary
operations can be continued at any precision.

We finally specify the scope of this calculation
for the requested further quotients. Since
$N''\subset\gamma_4N$, there is the actual arithmetic
surjection
\[
 N/N''\longrightarrow Q_3,
       \qquad\ker=\gamma_4N/N''.
 \tag{E40.13}
\]
For the single rational cusp marking $I_Q$, the image
of $I_Q$ in $V$ has rank one. Hence its commutator
subgroup maps to zero in
$\gamma_2N/\gamma_3N=\bigwedge^2V$, and
$[I_Q,I_Q]\subset\gamma_3N$.
It follows that $\gamma_3I_Q\subset\gamma_4N$, giving
\[
 N/\gamma_3I_Q\longrightarrow Q_3,
       \qquad\ker=\gamma_4N/\gamma_3I_Q.
 \tag{E40.14}
\]
The maps (E40.13)--(E40.14), followed by $Q_3\to Q_2$,
identify exactly the common stage whose fibers were
computed. They do not identify either full kernel
with the degree-three central lattice.
\end{proof}
